\scanpage{149}
\mspage{98}

moins que ceci : $f_!$ transforme les monomorphismes $V'\unsure{-|\to}V$, $V''\mathbin{-}V$ (interprétés comme flèches dans $V^\wedge$, de but $e_V$) en monomorphismes, et le produit fibré $V'\cap V''=V'\times_V V''$ en \inserted{le} produit fibré correspondant dans $W^\wedge$. Mais le foncteur $f_!$ commute aux $\varinjlim$ quelc. (étant un adjoint à g.), et je ne vois aucune raison a priori pour qu'il \deleted{\unsure{trans}} commute \inserted{aux} \inserted{\unsure{certains}} produits fibrés (ce qui impliquerait aussi qu'il transforme monos. en \unsure{mono|monos.}). Si je me rappelle bien, c'est ultra faux en général.

Donc je doute fort à présent que l'énoncé b) soit vrai, sans hyp. supplémentaires. Pour assurer la monomorphie, on songe à imposer l'existence de \uline{rétractions} de $V$ sur $V',V''$ et de $V',V''$ sur $V'\cap V''$.

Je \unsure{viens|tâche} de vérifier la chose, et donne le

\uline{Corollaire} (rectification de prop 9 partie b) : l'énoncé est valable si on suppose \deleted{\unsure{l'exi}} qu'il existe une \uline{rétraction} $r'_0$ de $V''$ sur $V'\cap V''$, et de $=$ une rétraction\unsure{$^{\mathrm{P!}}$|!} $r''_0$ de $V'$ sur $V'\cap V''$.

% slanted
\margin{NB Il y a un énoncé légèrement plus général, en partant d'une catégorie $U$ et d'un diagramme \diagram{\begin{tikzcd}[ampersand replacement=\&,column sep=small,row sep=small]
V'\cap V''\arrow[r]\arrow[d]\&V'\arrow[d]\\
V''\arrow[r]\&V
\end{tikzcd}} \deleted{\unsure{de}} $\unsure{U|U^\wedge}$ : Si $V',V'',V'\cap V''$ sont $\unsure{W|\mathcal W}$-asph. \unsure{rel.} $U$, il en est de $=$ de $V$.}

\diagram{\begin{tikzcd}[ampersand replacement=\&]
 V'\cap V'' \arrow[r,"i''_0"'] \arrow[d,"i'_0"] \& V' \arrow[d,"i'"] \arrow[l,dashed,bend right=35,"r''_0"']\\
 V'' \arrow[r,"i''"] \arrow[u,dashed,bend left=38,"r'_0"] \& V \arrow[u,dashed,bend right=35,"r'"] \arrow[l,dashed,bend left=32,"r''"]
\end{tikzcd}}

(NB Nous verrons plus tard qu'il suffit de l'existence d'\uline{une} de ces rétractions pour assurer la conclusion)

% slanted
\margin{\unsure{Cor.} Mais peut-être faux également -- les rétractions \unsure{effet|effets} \unsure{de} \unsure{ce terme} sur les flèches en les passant \unsure{à} \unsure{$V^\wedge$|$V$} \unsure{ou} en \unsure{les}.}

\uline{Dém}. Comme $V$ est somme amalgamée de $\deleted{\unsure{V}}V'$ et $V''$, \deleted{l'existence de} il existe une rétraction $r'$ de $V$ sur $V'$ (donnée par $\mathrm{id}_{V'}$ sur $V'$, par $i''_0r'_0$ sur $V''$) qui induit $r'_0$ sur $V''$, et de $=$ il existe
% archive sheet number: 149
